Las siguientes alternativas están valoradas según su cumplimiento de criterios, a su utilidad práctica, limitaciones de datos, costo computacional y calendario.

\subsection{Acciones posibles}

\begin{table}[H]
  \centering
  \small
  \caption{Acciones posibles.}
  \label{tab:acciones}
  \begin{tabularx}{\textwidth}{>{\raggedright\arraybackslash}p{0.28\textwidth}>{\raggedright\arraybackslash}X>{\raggedright\arraybackslash}X}
    \toprule
    Acción & Ventajas & Desventajas \\
    \midrule
    A. Despliegue del piloto asistido con el Modelo 5 & Ya tenemos un modelo que supera las metas. Dicho piloto permitiría observar el uso real & El Modelo 5 sigue siendo débil a extremos, por lo que requeriría mantener una revisión por parte de un experto, tener seguimiento de errores y protocolos de captura. \\
    \addlinespace
    B. Iterar modelado antes de desplegar & Podría mejorarse el modelo en los extremos y concluir EVA02. & Coste de cálculo, mejoras inciertas y un potencial retraso. \\
    \addlinespace
    C. Recolectar más casos extremos y reentrenar & Atacaría la escasez de ejemplos de este tipo y permitiría tener una validación más representativa. & Requeriría de una nueva visita, etiquetado, tiempo de especialistas y volver a la fase de Preparación de los datos. \\
    \bottomrule
  \end{tabularx}
\end{table}
